% ============================================================================
% Fichier  : partie1/P01_C01_audit_controle_auditologie.tex
% Titre    : Auditer, Contrôler, Donner une Assurance : du Métier à l'Auditologie
% Auteur   : MINKA MI NGUIDJOÏ Thierry Emmanuel
% Version  : 1.0 -- novembre 2026
% ============================================================================
% Sources : CRIPP 2017 de l'IIA (définition de l'audit interne, paraphrasée) ;
% Normes mondiales de l'IIA (2024) ; modèle des trois lignes (IIA, 2020, six
% principes, paraphrasés) ; ISO 19011:2018 (vocabulaire : audit, critères,
% preuves, constatations, conclusions ; paraphrasé) ; guide d'audit des SI du
% CHAI (2014), typologie ; INTOSAI GUID 5100 ; loi n° 2010/012, texte
% français officiel, art. 4 (9) et 13 (1), cités comme dans ACSI C12 ;
% article A0 (HAL hal-05701007). Dates de l'ISACA (1969, 1978) et de COBIT
% (1996, 2007, 2012, 2018-2019) : présentations publiques de l'ISACA.
% Faits LiZbiZ communs avec ACSI C01 : 2,4 millions de mouvements de stock
% annuels ; lettre de mission d'AUDIX CONSEIL (Douala + 4 agences sur 12).
% Corrigés : corriges/Corr_C01.tex.
% ============================================================================

\chapter[Auditer, Contrôler, Donner une Assurance]{Auditer, Contrôler, Donner une Assurance : du Métier à l'Auditologie}
\label{chap:audit-auditologie}
\chaptermark{Du métier à l'auditologie}

\epigraph{%
    « L'auditeur ne dit pas que tout va bien.\\
    Il dit ce qu'il sait, comment il le sait, et ce qu'il ignore. »%
}{Adage de la profession, reformulé}

\niveauChapitre{Fondamental}{%
    Maîtrise du Chapitre~\ref{chap:prologue} (impression, constatation,
    opinion~; réponse \NEK{}~; mesures $\mu$, $U$, $H$)%
}

\begin{boxobjectifs}
\begin{itemize}
    \item Distinguer audit, contrôle, inspection, conseil, certification et
          expertise, par la question posée, l'acteur et le livrable.
    \item Utiliser correctement le vocabulaire de l'audit~: critère, preuve,
          constatation, conclusion, opinion.
    \item Définir l'assurance raisonnable et l'assurance limitée, rédiger une
          conclusion dans chacune des deux formes, et expliquer pourquoi
          aucune assurance n'est absolue.
    \item Estimer, par un calcul simple, le risque de ne pas voir une
          anomalie quand on travaille sur un échantillon.
    \item Situer l'audit interne, l'audit externe et l'audit public, et
          appliquer le modèle des trois lignes à une organisation.
    \item Retracer les grandes étapes de l'audit informatique et de sa
          profession.
    \item Classer une mission selon plusieurs critères (objet, objectif,
          moment, niveau d'assurance) et choisir le type d'audit adapté à un
          problème.
    \item Présenter l'audit comme réduction de l'incertitude sur un système
          sociotechnique, et ce que cette lecture change pour la pratique.
\end{itemize}
\end{boxobjectifs}

\bigskip

% ============================================================================
\section*{Question directrice : qu'achète-t-on quand on paie un audit~?}
% ============================================================================

Après l'échec de la première mission (Chapitre~\ref{chap:prologue}),
M.~ETEME, directeur général de LiZbiZ, reçoit M.~BEYALA, chef de mission
d'AUDIX CONSEIL. Sa question est simple~: «~Si je vous paie, pouvez-vous me
garantir qu'il n'y a plus de fraude dans mes entrepôts~?~» M.~BEYALA répond
non. Il peut garantir une démarche, des preuves et une opinion motivée sur
ce qu'il aura examiné. Il ne peut pas garantir l'absence de toute fraude.

Cette réponse déçoit souvent les dirigeants. Elle est pourtant au cœur du
métier, et ce chapitre l'explique. Il définit ce qu'est un audit et le
distingue des activités voisines (section~\ref{sec:c01-notions}), fixe son
vocabulaire (section~\ref{sec:c01-vocabulaire}), dit ce qu'il garantit et ce
qu'il ne garantit pas (section~\ref{sec:c01-assurance}), présente ceux qui
auditent (section~\ref{sec:c01-acteurs}), retrace l'histoire de l'audit
informatique (section~\ref{sec:c01-histoire}), classe les missions
(section~\ref{sec:c01-types}), situe le cadre camerounais
(section~\ref{sec:c01-cadre}), et propose enfin de lire l'audit comme une
\emph{réduction de l'incertitude} (section~\ref{sec:c01-auditologie}).

\separateur

% ============================================================================
\section{Audit, Contrôle, Inspection, Conseil}
\label{sec:c01-notions}
% ============================================================================

\subsection{Six activités voisines}

Dans une organisation, plusieurs activités examinent le fonctionnement des
systèmes. Elles se ressemblent, emploient parfois les mêmes personnes et
les mêmes outils, mais elles ne posent pas la même question et ne
produisent pas le même livrable. Les confondre est la source de la plupart
des malentendus entre auditeurs et audités.

\begin{boxdef}
\termecle{Audit}\index{audit}~: examen méthodique, indépendant et documenté,
qui confronte une situation à des critères connus à l'avance, et conduit à
une opinion motivée sur cette situation, fondée sur des preuves.

\medskip
\termecle{Contrôle}\index{contrôle}~: dispositif ou action, intégré au
fonctionnement de l'organisation, qui vise à maîtriser une activité au fil
de l'eau (vérifier une saisie, approuver un paiement, rapprocher deux
comptes). Le contrôle est exercé par l'organisation elle-même.

\medskip
\termecle{Inspection}\index{inspection}~: vérification ponctuelle de la
conformité d'un objet ou d'une opération à une règle, souvent par une
autorité de tutelle, et souvent à visée de sanction.

\medskip
\termecle{Conseil}\index{conseil}~: accompagnement qui aide l'organisation à
concevoir ou à améliorer un dispositif. Le conseiller propose~; il ne donne
pas d'opinion indépendante sur ce qu'il a contribué à concevoir.

\medskip
\termecle{Certification}\index{certification}~: attestation, par un
organisme tiers accrédité, qu'un système de management est conforme aux
exigences d'une norme, sur un périmètre et pour une durée donnés.

\medskip
\termecle{Expertise}\index{expertise}~: avis technique donné par un
spécialiste sur une question précise, souvent à la demande d'un juge ou
d'une partie à un litige. L'expert répond à la question posée~; il ne
porte pas d'opinion d'ensemble sur un dispositif.
\end{boxdef}

\begin{tableaucours}{L{2.3cm}L{3.6cm}L{3.0cm}L{3.3cm}}{%
    \tableauHeader{Activité} &
    \tableauHeader{Question posée} &
    \tableauHeader{Qui l'exerce} &
    \tableauHeader{Livrable}
}
Audit & La situation est-elle conforme aux critères, et dans quelle mesure
le sait-on~? & Auditeur indépendant de l'activité & Rapport et opinion
motivée \\
Contrôle & L'opération en cours est-elle correcte~? & Les opérationnels,
la hiérarchie & Une opération validée ou rejetée \\
Inspection & La règle est-elle respectée ici et maintenant~? & Une autorité
de tutelle & Procès-verbal, sanction éventuelle \\
Conseil & Comment faire mieux~? & Consultant, expert & Recommandations,
dispositif conçu \\
Certification & Le système satisfait-il la norme~? & Organisme accrédité &
Certificat, pour une durée donnée \\
Expertise & Que s'est-il passé techniquement~? & Expert désigné & Rapport
d'expertise \\
\end{tableaucours}
\vspace{-0.8em}
\begin{center}{\small\itshape Tableau~1.1. Six activités d'examen à ne pas
confondre}\end{center}

Une même personne peut, au fil de sa carrière, exercer plusieurs de ces
activités. Elle ne peut pas les exercer \emph{sur le même objet}. Celui qui
a conseillé la conception d'un dispositif ne peut pas l'auditer ensuite~;
celui qui exécute un contrôle ne peut pas en évaluer l'efficacité.

\begin{boxscenario}{Opération LIZBIZ. Qui fait quoi~?}
\textbf{Faits.} Au printemps 2026, cinq examens portent sur le
système d'information de LiZbiZ.
\begin{itemize}
    \item Chaque soir, le chef d'agence de Garoua vérifie que les bons de
          sortie de l'entrepôt sont signés.
    \item Mme~NKENG, contrôleuse interne, revoit chaque mois un échantillon
          d'ajustements de stock.
    \item AUDIX CONSEIL conduit la seconde mission.
    \item Le cabinet NDONGO \& ASSOCIÉS certifie les comptes de l'exercice.
    \item Un éditeur aide le service informatique à paramétrer le nouveau
          module de facturation.
\end{itemize}

\textbf{Analyse} (confiance élevée). Le chef d'agence exerce un
\emph{contrôle} de premier niveau~; Mme~NKENG, un contrôle de second niveau
(surveillance)~; AUDIX CONSEIL, un \emph{audit} contractuel~; NDONGO \&
ASSOCIÉS, un audit légal des comptes, qui s'appuie en partie sur les
contrôles informatiques~; l'éditeur, du \emph{conseil}. Si l'éditeur était
ensuite chargé d'auditer le module qu'il a paramétré, son opinion ne
vaudrait rien~: il se jugerait lui-même.
\end{boxscenario}

\subsection{L'audit interne selon l'IIA}

L'Institut des auditeurs internes (\sigle{IIA}) définit l'audit interne
comme une activité indépendante et objective, qui donne à une organisation
une assurance sur le degré de maîtrise de ses opérations, lui apporte des
conseils pour les améliorer et contribue à créer de la valeur. Il y parvient
en évaluant, par une approche systématique et méthodique, ses processus de
gouvernance, de management des risques et de contrôle~\cite{iia2017cripp}.
Les Normes mondiales de l'IIA, publiées en 2024, reprennent cette mission
et l'organisent en principes et en exigences~\cite{iia2024gias}~; le
Chapitre~\ref{chap:normes} les présente.

Trois mots de la définition méritent l'attention.
\begin{itemize}
    \item \emph{Assurance}~: l'audit produit un degré de confiance, pas une
          garantie (section~\ref{sec:c01-assurance}).
    \item \emph{Conseil}~: l'audit interne peut aussi conseiller, à
          condition de ne pas auditer ensuite ce qu'il a conçu
          (Chapitre~\ref{chap:deontologie}).
    \item \emph{Systématique}~: l'audit suit une méthode, que la
          Partie~\ref{part:methode} enseigne.
\end{itemize}

\begin{boxalert}
\textbf{Violation typique.} La contrôleuse interne de LiZbiZ rédige la
procédure de désactivation des comptes, la met en place, puis est chargée
«~d'auditer~» son application.

\smallskip
\textbf{Règle~:} on n'audite pas ce qu'on a conçu ou ce qu'on exécute. Le
contrôle relève de l'organisation~; l'audit doit en être indépendant.
\end{boxalert}

% ============================================================================
\section{Le Vocabulaire de l'Audit}
\label{sec:c01-vocabulaire}
% ============================================================================

La norme ISO~19011, qui donne les lignes directrices de l'audit des
systèmes de management, fixe un vocabulaire que ce cours
adopte~\cite{iso19011_2018}. Elle définit l'audit comme un processus
méthodique, indépendant et documenté, qui permet d'obtenir des preuves
objectives et de les évaluer objectivement, pour déterminer dans quelle
mesure les critères d'audit sont satisfaits. Cinq notions s'enchaînent.

\begin{boxdef}
\termecle{Critère d'audit}\index{critère d'audit}~: exigence de référence à
laquelle on compare ce que l'on observe~: une procédure interne, une
norme, un article de loi, un contrat.

\medskip
\termecle{Preuve d'audit}\index{preuve d'audit}~: information vérifiable,
pertinente pour le critère, sur laquelle l'auditeur fonde son appréciation
(document, extraction, observation, réponse corroborée).

\medskip
\termecle{Constatation d'audit}\index{constatation}~: résultat de la
comparaison des preuves avec le critère~: conformité, non-conformité, ou,
dans ce cours, impossibilité de conclure (\NEK{}).

\medskip
\termecle{Conclusion d'audit}\index{conclusion d'audit}~: résultat de la
mission, tiré de l'ensemble des constatations, au regard des objectifs de
la mission.

\medskip
\termecle{Opinion d'audit}\index{opinion d'audit}~: forme que prend la
conclusion lorsqu'elle est exprimée avec un niveau d'assurance
(section~\ref{sec:c01-assurance}).
\end{boxdef}

\begin{center}
\begin{tikzpicture}[
    noeud/.style={draw=CouleurPrimaire, thick, rounded corners,
                  fill=CouleurPrimaire!8, minimum width=2.3cm,
                  minimum height=0.9cm, align=center, font=\small},
    fleche/.style={-latex, thick, color=CouleurPrimaire}]
\node[noeud] (crit) at (0,1.2) {Critère};
\node[noeud] (preu) at (0,-0.2) {Preuve};
\node[noeud] (cons) at (3.4,0.5) {Constatation\\ \scriptsize\OUI{} / \NON{} / \NEK{}};
\node[noeud] (conc) at (6.8,0.5) {Conclusion};
\node[noeud] (opin) at (10.0,0.5) {Opinion\\ \scriptsize avec assurance};
\draw[fleche] (crit) -- (cons);
\draw[fleche] (preu) -- (cons);
\draw[fleche] (cons) -- (conc);
\draw[fleche] (conc) -- (opin);
\end{tikzpicture}

\smallskip
{\small\itshape Figure~1.1. De la preuve à l'opinion}
\end{center}

L'enchaînement a une conséquence pratique~: \emph{aucune constatation sans
critère}. Une observation qui n'est rapportée à aucun critère reste une
description~; elle ne dit ni conforme ni non conforme. C'était le
deuxième défaut de la première mission de LiZbiZ.

\begin{boxscenario}{Opération LIZBIZ. Du critère à l'opinion}
\textbf{Critère.} Procédure interne~: le compte d'un agent parti est
désactivé sous 48~heures.

\textbf{Preuve.} Extraction des départs signée par les ressources
humaines, journal des désactivations, avec leurs empreintes (PT-14).

\textbf{Constatation.} Sur 31~départs, 19 comptes désactivés sous
48~heures, 7 plus tard, 5 encore actifs~: \NON{}.

\textbf{Conclusion.} Le contrôle de désactivation des comptes (D.2.1) n'est
pas efficace sur la période.

\textbf{Opinion.} Elle portera sur l'ensemble du domaine~D, après le test
de ses autres contrôles~; une constatation isolée ne fait pas une opinion
sur un domaine.
\end{boxscenario}

% ============================================================================
\section{Ce que l'Audit Garantit, et ce qu'il ne Garantit pas}
\label{sec:c01-assurance}
% ============================================================================

\subsection{Assurance raisonnable et assurance limitée}

\begin{boxdef}
\termecle{Assurance raisonnable}\index{assurance raisonnable}~: niveau élevé,
mais non absolu, de confiance que l'auditeur exprime sous une forme
positive («~nous estimons que le dispositif est efficace~»), après des
travaux étendus.

\medskip
\termecle{Assurance limitée}\index{assurance limitée}~: niveau de confiance
plus faible, exprimé sous une forme négative («~nous n'avons rien relevé
qui laisse penser que\ldots~»), après des travaux moins étendus, comme une
revue documentaire et des entretiens.
\end{boxdef}

\begin{tableaucours}{L{2.8cm}L{4.6cm}L{4.8cm}}{%
    \tableauHeader{} &
    \tableauHeader{Assurance raisonnable} &
    \tableauHeader{Assurance limitée}
}
Travaux & Tests étendus des contrôles, tests sur les données &
Entretiens, revue documentaire, procédures analytiques \\
Forme de l'opinion & Positive & Négative \\
Exemple de formulation & « À notre avis, les contrôles d'accès de
l'application ont fonctionné efficacement de janvier à avril~2026, au
regard de la procédure interne. » & « Nous n'avons rien relevé qui nous
conduise à penser que les contrôles d'accès n'ont pas fonctionné
efficacement de janvier à avril~2026. » \\
Coût et durée & Élevés & Modérés \\
Usage typique & Audit légal des comptes, audit de contrôle interne & Revue
de conformité, examen préliminaire \\
\end{tableaucours}
\vspace{-0.8em}
\begin{center}{\small\itshape Tableau~1.2. Deux niveaux d'assurance}\end{center}

\begin{boxalert}
\textbf{Violation typique.} Écrire « nous n'avons rien relevé
d'anormal » après avoir seulement interrogé le responsable informatique,
puis laisser le lecteur comprendre que les contrôles sont efficaces.

\smallskip
\textbf{Règle~:} la forme de l'opinion dit le niveau d'assurance. Une
opinion négative doit être lue comme telle, et le rapport doit dire quels
travaux l'ont fondée.
\end{boxalert}

\subsection{Pourquoi l'assurance n'est jamais absolue}

Quatre raisons tiennent à la nature même de l'audit.

\begin{enumerate}
    \item \textbf{L'auditeur échantillonne.} Il n'examine pas toutes les
          opérations~: une anomalie peut se trouver hors de l'échantillon.
    \item \textbf{Les contrôles peuvent être contournés}, notamment par
          collusion entre plusieurs personnes, ou par un dirigeant qui passe
          outre.
    \item \textbf{L'auditeur juge.} Le choix des critères, des tests et du
          seuil de signification repose sur un jugement professionnel.
    \item \textbf{L'auditeur ne voit pas tout.} Une partie de l'information
          est absente, cachée ou contradictoire. C'est là que la réponse
          \NEK{} prend sa place (section~\ref{sec:c01-auditologie}).
\end{enumerate}

\subsection{Le risque de ne pas voir~: un calcul}

La première raison se chiffre. Supposons qu'une proportion $t$ des
opérations d'une population soient anormales, et que l'auditeur en tire au
hasard un échantillon de $n$ opérations. Si la population est grande devant
l'échantillon, chaque opération tirée est anormale avec la probabilité $t$,
indépendamment des autres. La probabilité de ne tirer \emph{aucune}
anomalie, donc de ne rien voir, vaut
\[
    P(\text{aucune anomalie vue}) = (1 - t)^{n}.
\]

\begin{tableaucours}{cccc}{%
    \tableauHeader{Taux d'anomalies $t$} &
    \tableauHeader{$n = 25$} &
    \tableauHeader{$n = 100$} &
    \tableauHeader{$n = 230$}
}
5~\% & 0,28 & 0,006 & $< 10^{-5}$ \\
2~\% & 0,60 & 0,13 & 0,010 \\
1~\% & 0,78 & 0,37 & 0,099 \\
\end{tableaucours}
\vspace{-0.8em}
\begin{center}{\small\itshape Tableau~1.3. Probabilité de ne voir aucune
anomalie dans un échantillon aléatoire}\end{center}

Avec 25 opérations tirées et 2~\% d'anomalies, l'auditeur a six chances sur
dix de ne rien voir. Pour ramener ce risque à 10~\% avec 1~\% d'anomalies,
il faut $n$ tel que $0{,}99^{n} \leq 0{,}10$, soit
\[
    n \geq \frac{\ln 0{,}10}{\ln 0{,}99} \approx 229{,}1,
    \qquad\text{donc}\qquad n = 230.
\]
Ce calcul simple est le cœur de l'échantillonnage statistique par
attributs~; le Chapitre~\ref{chap:preuve} le développera (tailles
d'échantillon, niveaux de confiance, erreurs tolérables). Il montre déjà
pourquoi l'assurance ne peut pas être absolue~: pour être \emph{certain} de
voir une anomalie rare, il faudrait tout examiner.

\begin{boiteRemarque}{Pourquoi ne pas tout vérifier~?}
Parce que le coût serait démesuré, et que le résultat ne serait pas
meilleur~: vérifier les 2,4~millions de mouvements de stock annuels de
LiZbiZ ne dirait rien des mouvements qui n'ont \emph{pas} été enregistrés.
L'exhaustivité porte sur ce qui est enregistré~; l'audit porte aussi sur ce
qui manque. Les outils d'analyse de données permettent aujourd'hui
d'examiner des populations entières (Chapitre~\ref{chap:audit-continu}),
mais ils ne voient, eux aussi, que ce que le système a enregistré.
\end{boiteRemarque}

\subsection{L'écart d'attentes}

On appelle \termecle{écart d'attentes}\index{écart d'attentes}
(\textit{expectation gap}) la différence entre ce que le public, ou le
commanditaire, croit qu'un audit garantit, et ce qu'il garantit réellement.
La question de M.~ETEME en est l'exemple type~: il attend une garantie
d'absence de fraude, qu'aucun audit ne peut donner. Le seul remède est
d'écrire, avant la mission, ce qu'elle fera et ce qu'elle ne fera pas.

\begin{boxscenario}{Opération LIZBIZ. La lettre de mission}
\textbf{Faits.} Le directeur général demande une garantie d'absence de
fraude.

\medskip
\textbf{Réponse du chef de mission.} La lettre de mission proposée par
AUDIX CONSEIL précise~:
\begin{itemize}
    \item l'objectif~: donner une assurance raisonnable sur l'efficacité des
          contrôles d'accès et des contrôles de gestion des stocks de
          l'application intégrée, sur la période de janvier à avril~2026~;
    \item le périmètre~: le centre de données de Douala et quatre agences
          sur douze, choisies sur la base des écarts d'inventaire~;
    \item ce que la mission ne fera pas~: elle ne recherche pas
          systématiquement la fraude, et ne remplace pas une enquête si des
          indices sérieux apparaissent~;
    \item les obligations de LiZbiZ~: accès aux personnes, aux documents et
          aux systèmes, et désignation d'un interlocuteur.
\end{itemize}

\textbf{Analyse} (confiance élevée). La lettre transforme une attente
irréaliste (« plus de fraude ») en un objectif vérifiable. Elle protège
l'auditeur, et elle protège aussi le directeur général, qui sait ce qu'il
achète.
\end{boxscenario}

% ============================================================================
\section{Qui Audite~? Interne, Externe, Public}
\label{sec:c01-acteurs}
% ============================================================================

\subsection{Les familles d'auditeurs}

\begin{tableaucours}{L{2.6cm}L{3.8cm}L{2.9cm}L{2.9cm}}{%
    \tableauHeader{Famille} &
    \tableauHeader{Rôle} &
    \tableauHeader{Commanditaire} &
    \tableauHeader{Exemple}
}
Audit interne & Assurance et conseil sur la maîtrise des opérations &
Organe de gouvernance (comité d'audit) & Direction de l'audit interne de la
BANQUE SANAGA \\
Audit externe légal & Certification des comptes & Assemblée des associés,
la loi & Le commissaire aux comptes de LiZbiZ \\
Audit externe contractuel & Mission ponctuelle sur un sujet défini & La
direction & AUDIX CONSEIL chez LiZbiZ \\
Audit de certification & Attestation de conformité à une norme de
management & L'organisation qui veut être certifiée & Certification ISO/IEC
27001 d'un centre de données \\
Audit réglementaire de sécurité & Audit de sécurité imposé par la loi &
La loi, l'autorité de régulation & Audit de sécurité au titre de la loi
n°~2010/012 \\
Audit public & Contrôle de l'emploi des fonds publics & La loi, la
juridiction & La juridiction des comptes à l'ANPP \\
\end{tableaucours}
\vspace{-0.8em}
\begin{center}{\small\itshape Tableau~1.4. Les familles d'auditeurs}\end{center}

Dans le secteur public, l'audit externe relève des institutions supérieures
de contrôle des finances publiques, qui appliquent les normes de
l'INTOSAI. Le GUID~5100 de l'INTOSAI fournit le cadre général de la
vérification des systèmes d'information~\cite{intosai2019guid5100}. Au
Cameroun, la Chambre des comptes de la Cour suprême exerce le contrôle
juridictionnel des comptes publics~; le Contrôle Supérieur de l'État exerce
le contrôle administratif des services publics. Le
Chapitre~\ref{chap:droit} précise leurs compétences.

\subsection{Le modèle des trois lignes}

L'IIA a publié en 2020 le \termecle{modèle des trois lignes}\index{modèle des
trois lignes}~\cite{iia2020troislignes}. Il remplace un modèle plus ancien,
dit des «~trois lignes de défense~», dont il abandonne la métaphore
militaire~: les lignes ne sont plus des remparts successifs, mais des rôles
complémentaires. Le modèle repose sur six principes, que l'on peut
résumer ainsi~:
\begin{enumerate}
    \item la gouvernance exige des structures et des processus qui
          permettent de rendre compte, d'agir et de donner une assurance~;
    \item l'organe de gouvernance rend compte aux parties prenantes et
          supervise~;
    \item le management porte la maîtrise des risques~: la première ligne
          dans les activités, la deuxième en appui et en surveillance~;
    \item l'audit interne, troisième ligne, donne une assurance et des
          conseils indépendants et objectifs~;
    \item l'indépendance de la troisième ligne à l'égard du management est
          une condition de son objectivité~;
    \item tous les rôles concourent à créer et à protéger la valeur.
\end{enumerate}

\begin{tableaucours}{L{2.8cm}L{4.6cm}L{2.4cm}L{2.4cm}}{%
    \tableauHeader{Acteur} &
    \tableauHeader{Rôle} &
    \tableauHeader{LiZbiZ} &
    \tableauHeader{BANQUE SANAGA}
}
Organe de gouvernance & Fixe la direction, supervise, rend compte aux
parties prenantes & Conseil d'administration & Conseil, comité d'audit,
comité des risques \\
Première ligne & Porte et maîtrise les risques dans ses activités &
Responsable informatique, chefs d'agence & DSI, agences, opérations \\
Deuxième ligne & Expertise, appui, surveillance des risques & Contrôleuse
interne & Direction des risques, RSSI, conformité \\
Troisième ligne & Assurance et conseil indépendants & Audit interne (à
créer) & Direction de l'audit interne \\
Prestataires externes & Assurance complémentaire & Commissaire aux
comptes, AUDIX CONSEIL & Commissaire aux comptes, régulateur \\
\end{tableaucours}
\vspace{-0.8em}
\begin{center}{\small\itshape Tableau~1.5. Le modèle des trois lignes
appliqué aux deux cas privés}\end{center}

\begin{boiteRemarque}{LiZbiZ n'a pas d'audit interne~: est-ce grave~?}
Pas nécessairement~: une entreprise de taille moyenne peut confier la
troisième ligne à un cabinet mandaté. Mais l'affaire du prologue montre le
risque de la situation actuelle~: la première mission a été conduite par
une équipe qui dépendait de la direction qu'elle évaluait, sans méthode ni
indépendance. Le modèle ne prescrit pas une structure~; il prescrit des
rôles, et l'indépendance de celui qui donne l'assurance.
\end{boiteRemarque}

\subsection{Les métiers du cours dans ce paysage}

\begin{boxmetier}{Contrôleur interne}
\textbf{Se situer en deuxième ligne}
\begin{enumerate}
    \item Concevoir et documenter les contrôles avec la première ligne,
          sans les exécuter à sa place.
    \item Surveiller leur fonctionnement par des contrôles de second niveau
          (revue d'échantillons, indicateurs).
    \item Remonter les anomalies à la direction et au comité d'audit, sans
          filtre.
    \item Laisser l'audit, interne ou externe, évaluer le dispositif de façon
          indépendante~: le contrôleur interne ne s'audite pas lui-même.
\end{enumerate}
\end{boxmetier}

\begin{boxmetier}{Gestionnaire de risques}
\textbf{Deuxième ligne, côté risques}

Il tient la cartographie et le registre des risques, aide la première ligne
à les évaluer et suit les plans de traitement. Il est un interlocuteur
majeur de l'auditeur, qui part de sa cartographie pour orienter la
mission (Chapitre~\ref{chap:planifier}), mais il n'est pas auditeur~: il
est lui-même audité sur la qualité de son dispositif.
\end{boxmetier}

\begin{boxmetier}{RSSI}
\textbf{Entre deux lignes}

Le responsable de la sécurité des systèmes d'information est placé, selon
les organisations, sous la DSI (première ligne) ou sous la direction des
risques (deuxième ligne). Le second rattachement préserve mieux sa capacité
à contredire la DSI. Pour l'auditeur, la question n'est pas formelle~: il
vérifie que le RSSI peut effectivement alerter la direction.
\end{boxmetier}

\begin{boxmetier}{Auditeur de sécurité}
\textbf{Un auditeur, pas un administrateur}

L'auditeur de sécurité intervient en troisième ligne ou comme prestataire
externe. Il évalue la sécurité~; il ne la met pas en œuvre. S'il corrige
lui-même une configuration pendant l'audit, il cesse d'être auditeur pour
devenir administrateur, et son rapport perd sa valeur.
\end{boxmetier}

% ============================================================================
\section{D'où Vient l'Audit Informatique}
\label{sec:c01-histoire}
% ============================================================================

\subsection{Trois manières d'auditer avec l'ordinateur}

Lorsque l'informatique entre dans les services comptables, dans les années
1960, les auditeurs l'abordent d'abord comme une boîte noire. Trois
approches se succèdent, puis coexistent.

\begin{tableaucours}{L{3.2cm}L{5.0cm}L{4.0cm}}{%
    \tableauHeader{Approche} &
    \tableauHeader{Principe} &
    \tableauHeader{Limite}
}
Autour de l'ordinateur & On rapproche les entrées et les sorties sans
examiner les traitements & Aveugle aux traitements automatisés et aux
volumes \\
À travers l'ordinateur & On teste les contrôles programmés et la logique
des traitements & Demande des compétences techniques \\
Avec l'ordinateur & On utilise des outils pour analyser l'ensemble des
données (techniques d'audit assistées par ordinateur) & Demande des
extractions fiables \\
\end{tableaucours}
\vspace{-0.8em}
\begin{center}{\small\itshape Tableau~1.6. Les approches de l'audit
informatique}\end{center}

Un exemple les distingue. Pour s'assurer que l'application de LiZbiZ
calcule correctement les factures, l'auditeur «~autour de l'ordinateur~»
recalcule à la main quelques factures à partir des bons de livraison.
L'auditeur «~à travers l'ordinateur~» examine la table des tarifs et la
règle de calcul programmée, et vérifie qui peut les modifier. L'auditeur
«~avec l'ordinateur~» extrait les 18\,340 expéditions de l'année et les
rapproche toutes des factures par une requête~; le Chapitre~\ref{chap:audit-applications}
montrera qu'une telle requête révèle quatorze expéditions livrées et jamais
facturées.

\subsection{Une profession qui s'organise}

\begin{tableaucours}{cL{10.4cm}}{%
    \tableauHeader{Année} &
    \tableauHeader{Événement}
}
1969 & Fondation d'une association d'auditeurs du traitement électronique
des données, qui deviendra l'ISACA \\
1978 & Première certification d'auditeur des systèmes d'information
(\sigle{CISA}) \\
1996 & Première édition du référentiel COBIT \\
2002 & Loi Sarbanes-Oxley aux États-Unis~: les dirigeants des sociétés
cotées rendent compte de l'efficacité du contrôle interne sur
l'information financière~\cite{sox2002} \\
2007 & COBIT~4.1, très utilisé dans l'audit~\cite{itgi2007cobit41} \\
2012 & COBIT~5, qui distingue gouvernance et gestion \\
2018 & COBIT~2019 (publication en novembre 2018)~\cite{isaca2018cobitintro} \\
2020 & Modèle des trois lignes de l'IIA~; quatrième édition de l'ITAF de
l'ISACA~\cite{isaca2020itaf} \\
2024 & Normes mondiales de l'audit interne de l'IIA~\cite{iia2024gias} \\
\end{tableaucours}
\vspace{-0.8em}
\begin{center}{\small\itshape Tableau~1.7. Quelques jalons (d'après les
présentations publiques de l'ISACA et de l'IIA)}\end{center}

La loi Sarbanes-Oxley a joué un rôle particulier. En obligeant les
dirigeants à attester l'efficacité du contrôle interne, et les auditeurs à
l'évaluer, elle a fait des contrôles informatiques (accès, changements,
exploitation) un objet d'audit à part entière~: ce sont les
\emph{contrôles généraux informatiques}, présentés au
Chapitre~\ref{chap:controle-interne}. Le Chapitre~\ref{chap:defaillances}
raconte l'affaire qui a conduit à cette loi.

\subsection{Vers l'audit continu}

Avec les volumes de données actuels, l'audit ponctuel montre ses limites~:
une anomalie peut durer des mois entre deux missions. L'\termecle{audit
continu}\index{audit continu} analyse les données et les contrôles de façon
répétée, voire en permanence. Il ne remplace pas la mission~: il la nourrit
d'alertes et lui laisse le jugement. Le Chapitre~\ref{chap:audit-continu} y
est consacré.

% ============================================================================
\section{Classer les Missions}
\label{sec:c01-types}
% ============================================================================

\subsection{Quatre critères de classement}

Une mission d'audit se décrit selon plusieurs critères indépendants. Les
nommer avant de commencer évite de promettre ce que l'on ne fera pas.

\begin{tableaucours}{L{2.6cm}L{4.4cm}L{5.2cm}}{%
    \tableauHeader{Critère} &
    \tableauHeader{Modalités} &
    \tableauHeader{Exemple}
}
Objet & Gouvernance, application, projet, sécurité, production, continuité,
tiers & Les contrôles d'accès de l'application de LiZbiZ \\
Objectif & Conformité (la règle est-elle respectée~?), efficacité (le
dispositif atteint-il son but~?), performance (au meilleur coût~?) &
Efficacité des contrôles d'accès \\
Moment & A priori (avant la mise en œuvre), en cours, a posteriori &
A posteriori, sur janvier-avril~2026 \\
Niveau d'assurance & Raisonnable, limitée & Raisonnable \\
\end{tableaucours}
\vspace{-0.8em}
\begin{center}{\small\itshape Tableau~1.8. Quatre critères de classement
d'une mission}\end{center}

\subsection{Les types d'audit des SI selon leur objet}

Le guide d'audit des systèmes d'information du Comité d'harmonisation de
l'audit interne (\sigle{CHAI}) distingue plusieurs thèmes
d'audit~\cite{chai2014guide}. Le cours les reprend et les complète~; chacun
est traité dans la Partie~\ref{part:pratique}.

\begin{tableaucours}{L{4.0cm}L{6.2cm}c}{%
    \tableauHeader{Type d'audit} &
    \tableauHeader{Question centrale} &
    \tableauHeader{Chapitre}
}
Gouvernance et fonction informatique & La fonction informatique sert-elle les
objectifs de l'organisation~? & \ref{chap:audit-gouvernance} \\
Application en service et données & L'application produit-elle des données
exactes et complètes~? & \ref{chap:audit-applications} \\
Projets et marchés & Le projet livrera-t-il ce qui est attendu, dans les
délais et le budget~? & \ref{chap:audit-projets} \\
Système de management de la sécurité & La sécurité est-elle pilotée, et ce
pilotage fonctionne-t-il~? & \ref{chap:audit-smsi} \\
Sécurité technique & L'architecture et les configurations protègent-elles
le système~? & \ref{chap:audit-technique} \\
Test d'intrusion, audit de code & Un attaquant peut-il compromettre le
système~? & \ref{chap:intrusion} \\
Production et continuité & L'exploitation est-elle fiable, et le service
survivra-t-il à un sinistre~? & \ref{chap:audit-continuite} \\
Tiers et nuage & Les prestataires sont-ils maîtrisés~? &
\ref{chap:audit-tiers} \\
Conformité & Les obligations légales et réglementaires sont-elles
respectées~? & \ref{chap:audit-conformite} \\
\end{tableaucours}
\vspace{-0.8em}
\begin{center}{\small\itshape Tableau~1.9. Les types d'audit des SI
(d'après le guide du CHAI, complété)}\end{center}

\begin{boxscenario}{Opération LIZBIZ. Quel type d'audit~?}
\textbf{Faits.} Écart de stock de 186~millions FCFA, réserve du
commissaire aux comptes, réponses incohérentes sur le contrôle d'accès.

\medskip
\textbf{Analyse} (confiance élevée). La demande initiale («~tout le
système d'information~») n'est pas un type d'audit. La seconde mission
combine deux types~: un audit d'\emph{application en service}, centré sur
la gestion des stocks de l'application intégrée, et un audit de
\emph{sécurité technique}, limité au contrôle d'accès (domaine~D). Son
objectif est l'efficacité des contrôles, a posteriori, avec une assurance
raisonnable. Ce découpage détermine les référentiels, les compétences de
l'équipe et la durée.
\end{boxscenario}

% ============================================================================
\section{Le Cadre Camerounais en un Coup d'Œil}
\label{sec:c01-cadre}
% ============================================================================

Au Cameroun, plusieurs textes rendent l'audit ou le contrôle interne
obligatoires pour certaines organisations. Ils seront étudiés au
Chapitre~\ref{chap:droit}~; en voici les principaux.

\begin{itemize}
    \item La loi n°~2010/012 relative à la cybersécurité et à la
          cybercriminalité définit l'audit de sécurité et prévoit un audit
          obligatoire et périodique des systèmes d'information de certaines
          organisations~\cite{loi2010012}.
    \item Le règlement COBAC R-2016/04 fixe les exigences minimales du
          contrôle interne des établissements de crédit de la
          CEMAC~\cite{cobac2016r04}.
    \item Le droit OHADA des sociétés commerciales organise la mission du
          commissaire aux comptes et, pour certaines sociétés, le comité
          d'audit~\cite{ohada2014auscgie}.
    \item La loi n°~2024/017 sur la protection des données à caractère
          personnel crée des obligations que l'audit de conformité vérifie
          (Chapitre~\ref{chap:audit-conformite}).
\end{itemize}

\begin{boxjuris}
\textbf{\loicam{n°~2010/012}, article~4, point~9}~: \textit{« Audit de
sécurité~: examen méthodique des composantes et des acteurs de la sécurité,
de la politique, des mesures, des solutions, des procédures et des moyens
mis en oeuvre par une organisation, pour sécuriser son environnement,
effectuer des contrôles de conformité, des contrôles d'évaluation de
l'adéquation des moyens (organisationnels, techniques, humains,
financiers) investis au regard des risques encourus, d'optimisation, de
rationalité et de performance. »}

\medskip
\textbf{Article~13, alinéa~1}~: \textit{« Sont soumis à un audit de sécurité
obligatoire, les réseaux de communications électroniques et les systèmes
d'information, des opérateurs, des autorités de certification et des
fournisseurs de services de communications électroniques. »}

\medskip
\textbf{Portée pour l'audit.} La définition légale couvre trois des
objectifs du Tableau~1.8~: la conformité, l'efficacité (adéquation des
moyens aux risques) et la performance (optimisation, rationalité). Un
audit de sécurité au sens de la loi n'est donc pas une simple liste de
contrôles~: il apprécie aussi si les moyens sont proportionnés aux
risques.
\end{boxjuris}

% ============================================================================
\section{L'Audit comme Réduction de l'Incertitude}
\label{sec:c01-auditologie}
% ============================================================================

\subsection{Un constat~: un métier sans socle scientifique}

L'audit des SI s'apprend par compagnonnage, se codifie dans des
référentiels professionnels et se pratique par jugement d'expert. Ces
pratiques sont précieuses, mais, comme le soutient l'article fondateur de
ce cours~\cite{minka2026auditologie}, il manque à la discipline un socle
scientifique unificateur~: des axiomes explicites, des propositions
réfutables et des mesures communes pour son objet central.

\subsection{Une proposition~: l'auditologie}

\begin{boxdef}
\termecle{Mission auditologique}\index{mission auditologique}~: ensemble
formé par cinq éléments~: le \emph{système audité}~$\mathcal{S}$ (personnes,
technologies, flux d'information, règles)~; l'\emph{auditeur}~$\mathcal{A}$,
ou l'équipe d'audit, avec ses limites de traitement, de mémoire et
d'attention~; le \emph{canal d'information}~$\mathcal{I}$ par lequel les
preuves parviennent à l'auditeur~; la \emph{structure stratégique}~$\mathcal{G}$
des relations entre auditeur, audité et tiers~; la \emph{trajectoire}~$\mathcal{T}$
du système au fil des missions.

\medskip
\termecle{Auditologie}\index{auditologie}~: science appliquée de la
réduction de l'incertitude sur les missions
auditologiques~\cite{minka2026auditologie}.
\end{boxdef}

\begin{boxscenario}{Opération LIZBIZ. Les cinq éléments}
\begin{tableaucours}{cL{3.6cm}L{7.6cm}}{%
    \tableauHeader{} &
    \tableauHeader{Élément} &
    \tableauHeader{Seconde mission de LiZbiZ}
}
$\mathcal{S}$ & Système audité & L'application intégrée, ses utilisateurs
des douze agences, ses quatre centres de données, ses procédures \\
$\mathcal{A}$ & Auditeur & L'équipe d'AUDIX CONSEIL, ses compétences et ses
limites (temps, attention) \\
$\mathcal{I}$ & Canal d'information & Entretiens, extractions, journaux,
documents~; certains sont absents ou contradictoires \\
$\mathcal{G}$ & Structure stratégique & Les intérêts du DG, du responsable
informatique, des chefs d'agence, du commissaire aux comptes \\
$\mathcal{T}$ & Trajectoire & La première mission, la seconde, le cycle de
novembre~2026 \\
\end{tableaucours}
\end{boxscenario}

L'article pose six axiomes ontologiques, présentés en détail au
Chapitre~\ref{chap:auditologie}. Trois d'entre eux éclairent dès maintenant
la pratique.

\begin{itemize}
    \item[] \textbf{Axiome~1 (Observabilité partielle).} L'auditeur n'a accès
          qu'à une partie des états pertinents du système audité, et cette
          partie est elle-même partiellement opaque.
    \item[] \textbf{Axiome~4 (Interaction stratégique).} L'auditeur et l'audité
          poursuivent des intérêts distincts~; leurs actions se comprennent
          comme des coups dans un jeu à information imparfaite.
    \item[] \textbf{Axiome~6 (Opposabilité).} Toute conclusion d'une mission
          doit pouvoir être défendue devant un tiers qui a autorité pour la
          contester.
\end{itemize}

{\small\itshape Axiomes reformulés d'après~\cite{minka2026auditologie}, avec
la numérotation de l'article~; les six axiomes figurent au
Chapitre~\ref{chap:auditologie}.}

\subsection{Mesurer ce que l'audit réduit}

Si l'audit réduit l'incertitude, on doit pouvoir mesurer cette réduction.
Le prologue en a donné l'exemple. Sur le sous-domaine D.2 (cycle de vie
des comptes), les quatre interlocuteurs avaient répondu 1~\OUI{}, 1~\NON{}
et 2~\NEK{}~: $U = 0{,}50$, $H = 1{,}5$~bit. Le test du contrôle D.2.1 a
établi un \NON{} sur la population entière. Sur ce contrôle, l'auditeur ne
dépend plus des déclarations~: l'incertitude est levée, et c'est une
non-conformité \emph{connue} qui la remplace.

C'est ce que l'on attend d'un test~: qu'il transforme un \NEK{} en \OUI{}
ou en \NON{} prouvés. Un test qui laisse l'incertitude entière n'a rien
appris. Le Chapitre~\ref{chap:preuve} en tire une règle de conduite~:
diriger les tests là où $U$ est élevé, comme M.~BEYALA l'a fait sur D.2,
D.4 et D.5.

\subsection{Ce que cette lecture change pour la pratique}

\begin{tableaucours}{L{4.4cm}L{7.8cm}}{%
    \tableauHeader{Lecture classique} &
    \tableauHeader{Lecture auditologique}
}
L'audit vérifie la conformité & L'audit réduit l'incertitude~; la
conformité à un référentiel est un \emph{choix d'objectif}, pas un
fondement \\
Une question appelle OUI ou NON & L'observabilité partielle impose une
troisième réponse, \NEK{} (Chapitre~\ref{chap:preuve}) \\
L'indépendance est une règle déontologique & L'indépendance se déduit de
l'interaction stratégique et de l'opposabilité
(Chapitre~\ref{chap:deontologie}) \\
Le taux de conformité résume un domaine & Un domaine se lit par trois
mesures, $\mu$, $U$ et $H$ (Chapitre~\ref{chap:conclure}) \\
Une grille longue est une grille complète & La grille doit respecter les
limites cognitives de l'auditeur~: on décompose
(Chapitre~\ref{chap:decomposition}) \\
Une mission est un événement isolé & Un système se suit dans le temps~;
on compare les cycles (Chapitre~\ref{chap:suivi}) \\
\end{tableaucours}
\vspace{-0.8em}
\begin{center}{\small\itshape Tableau~1.10. Deux lectures de
l'audit}\end{center}

\begin{boiteRemarque}{L'auditologie remplace-t-elle les normes professionnelles~?}
Non. Elle ne prescrit pas le «~bon audit~» et n'invalide pas les cadres
existants~\cite{minka2026auditologie}. Les normes de l'IIA, de l'ISACA ou de
l'INTOSAI disent ce que l'auditeur doit faire~; l'auditologie cherche à
expliquer \emph{pourquoi} les méthodes fonctionnent ou échouent, et à le
vérifier par l'expérience. Ce cours enseigne les deux.
\end{boiteRemarque}

\begin{boxinfo}
\textbf{Ce que dit la recherche.} L'article \emph{Vers l'auditologie} (A0)
propose de traiter l'audit comme un objet scientifique~: des axiomes
explicites, des propositions qui peuvent être réfutées par l'expérience,
des mesures communes. Il range explicitement sa démarche dans la science du
design, qui construit des artefacts (ici, des instruments et des méthodes
d'audit) et les évalue. L'article sur l'incertitude épistémique (A1) en
donne le premier instrument mesurable. La validation empirique de ces
propositions est en cours~; le cours dit, à chaque fois, ce qui est établi
et ce qui reste à valider (Chapitre~\ref{chap:incertitude}).
\end{boxinfo}

\begin{boxmetier}{Auditeur SI}
\textbf{Appliquer la lecture auditologique dès la première réunion}
\begin{enumerate}
    \item Annoncer à l'audité que la réponse \NEK{} est admise, et qu'elle
          sera documentée comme telle.
    \item Distinguer, dans le plan de travail, ce qui sera prouvé, ce qui
          sera seulement déclaré, et ce qui restera inconnu.
    \item Prévoir, pour chaque domaine, le calcul de $\mu$, $U$ et $H$.
    \item Rédiger chaque conclusion de façon qu'un tiers puisse la refaire à
          partir des papiers de travail.
\end{enumerate}
\end{boxmetier}

\begin{boxmetier}{Responsable de la conformité}
\textbf{Ce que l'audit lui apporte}

Le responsable de la conformité (et, pour les données personnelles, le
délégué à la protection des données) doit pouvoir démontrer aux autorités
que les obligations sont respectées. L'audit lui fournit cette
démonstration, à condition qu'elle soit opposable~: des critères tirés des
textes, des preuves conservées, des \NEK{} déclarés plutôt que dissimulés.
\end{boxmetier}

\separateur

\begin{boxresume}
\textbf{Ce qu'il faut retenir de ce chapitre~:}
\begin{itemize}
    \item \textbf{Audit, contrôle, inspection, conseil, certification,
          expertise} sont six activités distinctes~; l'audit est
          indépendant et conclut sur des preuves au regard de critères.
    \item \textbf{Le vocabulaire}~: critère et preuve donnent une
          constatation~; les constatations donnent une conclusion~; la
          conclusion exprimée avec un niveau d'assurance est une opinion.
          Pas de constatation sans critère.
    \item \textbf{L'assurance n'est jamais absolue}~: l'auditeur
          échantillonne (probabilité de ne rien voir $(1-t)^n$), juge, ne
          voit pas tout, et les contrôles peuvent être contournés.
    \item \textbf{Assurance raisonnable} (opinion positive, travaux étendus)
          ou \textbf{limitée} (opinion négative, travaux restreints)~; la
          lettre de mission réduit l'écart d'attentes.
    \item \textbf{Familles d'auditeurs}~: interne, externe légal ou
          contractuel, de certification, réglementaire, public~; et
          \textbf{trois lignes} de maîtrise des risques sous l'organe de
          gouvernance.
    \item \textbf{L'audit informatique} est passé de l'audit autour de
          l'ordinateur à l'audit avec l'ordinateur, puis à l'audit continu~;
          la loi Sarbanes-Oxley a fait des contrôles informatiques un objet
          d'audit.
    \item \textbf{Une mission se classe} par son objet, son objectif, son
          moment et son niveau d'assurance.
    \item \textbf{L'auditologie} lit l'audit comme une réduction de
          l'incertitude sur un système partiellement observable, peuplé
          d'acteurs stratégiques~; un bon test transforme un \NEK{} en
          \OUI{} ou en \NON{} prouvés.
\end{itemize}
\end{boxresume}

\begin{boxinfo}
\textbf{Points de contrôle.} Vous avez compris ce chapitre si vous savez~:
\begin{itemize}
    \item dire, pour une situation donnée, de laquelle des six activités
          elle relève~;
    \item rédiger une même conclusion sous forme d'assurance raisonnable et
          d'assurance limitée~;
    \item calculer la probabilité de ne voir aucune anomalie dans un
          échantillon, et la taille d'échantillon qui la ramène sous un
          seuil~;
    \item placer les acteurs d'une organisation dans le modèle des trois
          lignes~;
    \item classer une mission selon ses quatre critères~;
    \item expliquer pourquoi un test est utile quand $U$ est élevé.
\end{itemize}
\end{boxinfo}

% ============================================================================
\section*{Exercices}
% ============================================================================

\begin{boxexo}[title={\ding{45}\quad Exercice~1.1. Quelle activité~? \niveaubase}]
Pour chacune des situations suivantes, indiquez s'il s'agit d'un audit,
d'un contrôle, d'une inspection, d'un conseil, d'une certification ou d'une
expertise, et justifiez.
\begin{enumerate}[(a)]
    \item Le chef d'agence de Garoua vérifie chaque soir que les bons de
          sortie de l'entrepôt sont signés.
    \item La COBAC se rend dans une agence de la BANQUE SANAGA pour vérifier
          le respect d'une règle de conservation des pièces.
    \item AUDIX CONSEIL donne une opinion sur l'efficacité des contrôles
          d'accès de LiZbiZ.
    \item Un cabinet aide l'ANPP à rédiger sa politique de sécurité.
    \item Un organisme accrédité délivre à un centre de données un
          certificat ISO/IEC~27001 valable trois ans.
    \item Un tribunal désigne un spécialiste pour dire comment un fichier a
          été modifié.
\end{enumerate}
\end{boxexo}

\begin{boxexo}[title={\ding{45}\quad Exercice~1.2. Critère, preuve, constatation, conclusion \niveaubase}]
Classez chacun des éléments suivants, tirés du dossier de la seconde
mission de LiZbiZ, dans l'une des quatre catégories, et justifiez.
\begin{enumerate}[(a)]
    \item « La revue des comptes d'administration est réalisée chaque
          trimestre par le responsable de la sécurité. » (procédure interne)
    \item L'extraction des comptes d'administration du 15~mai~2026, avec son
          empreinte SHA-256.
    \item « Aucun compte rendu de revue n'a été produit pour le premier
          trimestre~2026. »
    \item « Le contrôle des comptes à privilèges n'est pas efficace sur la
          période. »
    \item Le journal des connexions de l'application, de janvier à
          avril~2026.
\end{enumerate}
\end{boxexo}

\begin{boxexo}[title={\ding{45}\quad Exercice~1.3. Le risque de ne pas voir \niveaumoyen}]
Une population de 6\,000 bons de sortie contient une proportion inconnue
$t$ de bons non signés.
\begin{enumerate}[(a)]
    \item Calculez la probabilité de ne voir aucun bon non signé dans un
          échantillon aléatoire de 50 bons, si $t = 3~\%$.
    \item Quelle taille d'échantillon ramène cette probabilité sous 5~\%,
          pour $t = 3~\%$~?
    \item L'auditeur tire 50 bons et n'en trouve aucun non signé. Peut-il
          conclure que tous les bons sont signés~? Que peut-il dire~?
    \item Pourquoi le calcul suppose-t-il que la population est grande
          devant l'échantillon~? Que se passe-t-il si l'on tire 50 bons sur
          60~?
\end{enumerate}
\end{boxexo}

\begin{boxexo}[title={\ding{45}\quad Exercice~1.4. Raisonnable ou limitée \niveaumoyen}]
Après une semaine d'entretiens et de revue documentaire sur la sauvegarde
des centres de données de LiZbiZ, sans aucun test de restauration,
l'équipe propose la conclusion~: « Les sauvegardes sont efficaces. »
\begin{enumerate}[(a)]
    \item Quel niveau d'assurance les travaux permettent-ils~? Justifiez.
    \item Réécrivez la conclusion dans la forme qui convient.
    \item Quels travaux faudrait-il ajouter pour pouvoir donner une
          assurance raisonnable~?
\end{enumerate}
\end{boxexo}

\begin{boxexo}[title={\ding{45}\quad Exercice~1.5. Les trois lignes de la BANQUE SANAGA \niveaumoyen}]
La BANQUE SANAGA compte une DSI, un RSSI rattaché au directeur des
risques, une direction de l'audit interne dirigée par Mme~EKAMBI, un
comité d'audit et un commissaire aux comptes.
\begin{enumerate}[(a)]
    \item Placez chacun de ces acteurs dans le modèle des trois lignes.
    \item Mme~EKAMBI a participé, en 2025, à la rédaction de la procédure de
          gestion des changements. Quel problème cela pose-t-il pour un
          audit de cette procédure~?
    \item Proposez une solution conforme à l'indépendance de l'audit.
    \item Le directeur général propose de rattacher le RSSI à la DSI « pour
          simplifier ». Quel est le risque pour la maîtrise de la
          sécurité~?
\end{enumerate}
\end{boxexo}

\begin{boxexo}[title={\ding{45}\quad Exercice~1.6. Classer cinq missions \niveaumoyen}]
Classez chacune des missions suivantes selon les quatre critères du
Tableau~1.8, et indiquez son type selon le Tableau~1.9.
\begin{enumerate}[(a)]
    \item Avant le lancement du nouveau système central de la BANQUE
          SANAGA, vérifier que la recette a été conduite selon la procédure.
    \item À la demande du ministère de tutelle, vérifier que l'ANPP a
          appliqué le code des marchés publics pour l'achat de ses serveurs.
    \item Pour LiZbiZ, dire si le prestataire d'hébergement respecte les
          clauses de son contrat.
    \item Pour un opérateur de télécommunications, réaliser l'audit de
          sécurité prévu par la loi n°~2010/012.
    \item Pour la BANQUE SANAGA, dire si le coût de l'exploitation
          informatique est maîtrisé.
\end{enumerate}
\end{boxexo}

\begin{boxexo}[title={\ding{45}\quad Exercice~1.7. Approches de l'audit informatique \niveaumoyen}]
On veut s'assurer que l'application de l'ANPP ne paie aucun agent après la
date de son départ.
\begin{enumerate}[(a)]
    \item Décrivez ce que ferait un auditeur «~autour de l'ordinateur~»,
          «~à travers l'ordinateur~» et «~avec l'ordinateur~».
    \item Quelle approche donne la preuve la plus forte~? À quelle
          condition~?
    \item Écrivez, en pseudo-code ou en SQL, la requête que lancerait
          l'auditeur «~avec l'ordinateur~», à partir d'une table des paiements
          et d'une table des agents.
\end{enumerate}
\end{boxexo}

\begin{boxexo}[title={\ding{45}\quad Exercice~1.8. Deux lectures de l'audit \niveaumoyen}]
Pour chacun des énoncés suivants, dites s'il relève de la lecture
classique ou de la lecture auditologique, et ce qu'il changerait à la
mission.
\begin{enumerate}[(a)]
    \item « Il faut au moins 80~\% de réponses positives pour déclarer un
          domaine conforme. »
    \item « Un domaine où un quart des réponses sont \NEK{} ne peut pas être
          déclaré conforme, quel que soit le taux de \OUI{}. »
    \item « Notre grille de 300 questions couvre tout le système. »
    \item « Nous reviendrons dans six mois pour voir si la situation a
          changé, et nous comparerons. »
\end{enumerate}
\end{boxexo}

\begin{boxexo}[title={\ding{45}\quad Exercice~1.9. Une lettre de mission publique \niveauexpert}]
La juridiction des comptes annonce un audit des systèmes d'information de
l'ANPP.
\begin{enumerate}[(a)]
    \item Rédigez l'objectif, le périmètre et les limites de la mission, sur
          le modèle de la lettre d'AUDIX CONSEIL.
    \item Indiquez le niveau d'assurance visé (raisonnable ou limitée) et
          justifiez.
    \item Identifiez deux situations dans lesquelles l'équipe devra répondre
          \NEK{}, et expliquez ce que ces réponses apprendront à la
          juridiction.
    \item Montrez en quoi l'axiome d'opposabilité s'applique avec une force
          particulière à un audit public.
\end{enumerate}
\end{boxexo}

\begin{boxexo}[title={\ding{45}\quad Exercice~1.10. De la loi aux critères \niveauexpert}]
Partez de la définition légale de l'audit de sécurité (loi n°~2010/012,
article~4, point~9).
\begin{enumerate}[(a)]
    \item Repérez les objets que la définition demande d'examiner, et les
          trois natures de contrôle qu'elle mentionne.
    \item Pour un centre de données de LiZbiZ, déduisez-en trois critères
          d'audit, un par nature de contrôle, et formulez pour chacun une
          question atomique cotable \OUI{}, \NON{} ou \NEK{}.
    \item Lequel des trois critères est le plus difficile à rendre
          objectif~? Pourquoi, et comment le rendre vérifiable~?
\end{enumerate}
\end{boxexo}

\bigskip

\begin{boxinfo}
\textbf{Pour aller plus loin.}
\begin{itemize}
    \item La norme ISO~19011~\cite{iso19011_2018}, articles~3 et~4~: le
          vocabulaire et les principes de l'audit.
    \item Le modèle des trois lignes de l'IIA~\cite{iia2020troislignes}~:
          un document court, qui se lit en une heure.
    \item Le guide d'audit des SI du CHAI~\cite{chai2014guide}, pour la
          typologie des missions et leurs fiches.
    \item L'article \emph{Vers l'auditologie}~\cite{minka2026auditologie},
          sections~1 à~3~: le constat, la définition et les axiomes.
\end{itemize}
\end{boxinfo}

\begin{boxinfo}
\textbf{Transition.} L'audit évalue un dispositif~: le contrôle interne.
Avant d'apprendre à l'auditer, il faut voir ce qui arrive quand il manque.
Le Chapitre~\ref{chap:defaillances} raconte quatre grandes défaillances, et
ce qu'elles ont changé dans la profession.
\end{boxinfo}
